\documentclass[10pt,a4paper]{amsart}
\usepackage{amsmath,amssymb,amsthm,graphicx}
% Retained minimal preamble for the Stacks Project chapter more-algebra.tex (More on Algebra).
% The chapter file pinned-original-source.tex carries only \input{preamble} on its first line, so no
% preamble text lives inside the chapter itself.  The Stacks Project uses one single shared preamble for
% every chapter; the snapshot stages it identically as sources/stacks-chow/preamble.tex and
% sources/stacks-etale-cohomology/preamble.tex.  The block below is copied mechanically, line for line,
% from sources/stacks-chow/preamble.tex lines 14-18 (the Xy-pic setup) and 154-217 (the native theorem
% declarations together with the mathematical macros).  Nothing was added, removed or redefined; the
% wrapper class and the amsmath/amssymb/amsthm/graphicx base packages are supplied by the assembler.
% Chapter: more-algebra.tex -- More on Algebra.
% For commutative diagrams we use Xy-pic
\usepackage[all]{xy}

% We use 2cell for 2-commutative diagrams.
\xyoption{2cell}

% Theorem environments.
%
\theoremstyle{plain}
\newtheorem{theorem}[subsection]{Theorem}
\newtheorem{proposition}[subsection]{Proposition}
\newtheorem{lemma}[subsection]{Lemma}

\theoremstyle{definition}
\newtheorem{definition}[subsection]{Definition}
\newtheorem{example}[subsection]{Example}
\newtheorem{exercise}[subsection]{Exercise}
\newtheorem{situation}[subsection]{Situation}

\theoremstyle{remark}
\newtheorem{remark}[subsection]{Remark}
\newtheorem{remarks}[subsection]{Remarks}

\numberwithin{equation}{subsection}

% Macros
%
\def\lim{\mathop{\mathrm{lim}}\nolimits}
\def\colim{\mathop{\mathrm{colim}}\nolimits}
\def\Spec{\mathop{\mathrm{Spec}}}
\def\Hom{\mathop{\mathrm{Hom}}\nolimits}
\def\Ext{\mathop{\mathrm{Ext}}\nolimits}
\def\SheafHom{\mathop{\mathcal{H}\!\mathit{om}}\nolimits}
\def\SheafExt{\mathop{\mathcal{E}\!\mathit{xt}}\nolimits}
\def\Sch{\mathit{Sch}}
\def\Mor{\mathop{\mathrm{Mor}}\nolimits}
\def\Ob{\mathop{\mathrm{Ob}}\nolimits}
\def\Sh{\mathop{\mathit{Sh}}\nolimits}
\def\NL{\mathop{N\!L}\nolimits}
\def\CH{\mathop{\mathrm{CH}}\nolimits}
\def\proetale{{pro\text{-}\acute{e}tale}}
\def\etale{{\acute{e}tale}}
\def\QCoh{\mathit{QCoh}}
\def\Ker{\mathop{\mathrm{Ker}}}
\def\Im{\mathop{\mathrm{Im}}}
\def\Coker{\mathop{\mathrm{Coker}}}
\def\Coim{\mathop{\mathrm{Coim}}}

% Boxtimes
%
\DeclareMathSymbol{\boxtimes}{\mathbin}{AMSa}{"02}

%
% Macros for moduli stacks/spaces
%
\def\QCohstack{\mathcal{QC}\!\mathit{oh}}
\def\Cohstack{\mathcal{C}\!\mathit{oh}}
\def\Spacesstack{\mathcal{S}\!\mathit{paces}}
\def\Quotfunctor{\mathrm{Quot}}
\def\Hilbfunctor{\mathrm{Hilb}}
\def\Curvesstack{\mathcal{C}\!\mathit{urves}}
\def\Polarizedstack{\mathcal{P}\!\mathit{olarized}}
\def\Complexesstack{\mathcal{C}\!\mathit{omplexes}}
% \Pic is the operator that assigns to X its picard group, usage \Pic(X)
% \Picardstack_{X/B} denotes the Picard stack of X over B
% \Picardfunctor_{X/B} denotes the Picard functor of X over B
\def\Pic{\mathop{\mathrm{Pic}}\nolimits}
\def\Picardstack{\mathcal{P}\!\mathit{ic}}
\def\Picardfunctor{\mathrm{Pic}}
\def\Deformationcategory{\mathcal{D}\!\mathit{ef}}

% --- appended 2026-10-12 (Chief): the project own definitions of the reference/slogan/history
% display environments, copied mechanically from the shared Stacks preamble lines 8-12
% (\usepackage{verbatim} plus the three \newenvironment lines, which the chapter preamble keeps
% verbatim from the project preamble). Nothing added, removed or redefined: these environments
% carry no mathematical content and the project itself typesets them as comments. Without these
% three lines a counted statement whose own text uses them fails with Environment slogan undefined.
\usepackage{verbatim}
\newenvironment{reference}{\comment}{\endcomment}
\newenvironment{slogan}{\comment}{\endcomment}
\newenvironment{history}{\comment}{\endcomment}
\title{Direct summands of sums of quotients of the ring by principal ideals, characterized by a lifting property for short exact sequences}
\author{The Stacks Project Authors}
\date{stacks-project:more-algebra.tex:local-snapshot-20261012; GNU FDL 1.2 or later}
\begin{document}
\maketitle
\setcounter{subsection}{1}
\noindent\textit{Editorial note.} All delivered mathematical text is the author's own, copied from the pinned Stacks Project chapter file more-algebra.tex; no statement or proof step was rewritten, reordered or removed. The counted statement is the labeled lemma environment at source lines 37585-37598 (canonical tag 0ASM) and the counted proof is the proof environment at source lines 37600-37635, which immediately follows it and is its own proof.\par\medskip
\noindent\textit{Editorial note.} The counted statement contains the author's own reference environment, whose body is a citation of the paper by Warfield; the retained preamble defines the reference, slogan and history environments as comment-type (verbatim plus comment), so that body is discarded when the document is typeset and the citation is never evaluated. This is why the card needs no bibliography, and the environment is kept verbatim in the counted span as the author wrote it.\par\medskip
\noindent\textit{Editorial note.} The counted span contains no cross reference to any labelled item, so there is no navigation stub; the six references listed in other cards of this section do not occur here.\par\medskip
\noindent\textit{Editorial note.} The preamble span runs to line 90, the current end of the retained preamble, which is exactly what makes the reference environment a comment; the excerpt serves only as preamble, never as a span coordinate.\par\medskip
\noindent\textit{Editorial note.} The source URL is the source-level Stacks Project URL, as required when more than one card is produced from one source; the per-claim tag is carried by claim\_key and by this note.\par\medskip
\noindent\textit{Allowed standard prerequisites (editorial).}\begin{itemize}
\item The hypothesis that the intersection of a submodule with the multiples of an element equals the multiples of the submodule is what makes the lifting argument for the quotients work. used at proof lines 37601-37607. Proves the first implication by an explicit lift constructed from the divisibility of the image of one.
\item The canonical presentation of the module as the cokernel of the map from the direct sum over the pairs of an element and a map out of the corresponding quotient. used at proof lines 37610-37622. Produces the short exact sequence to which the second condition is applied in the converse.
\item A relation in the canonical presentation can be corrected by an explicit class in a single summand, making the corrected tuple map to the original element. used at proof lines 37622-37634. Completes the converse by exhibiting the element whose image is the given one.
\end{itemize}\medskip
\begin{lemma}
\label{lemma-characterize-PD-modules}
\begin{reference}
\cite[Corollary 1]{Warfield-Purity}
\end{reference}
Let $P$ be a module over a ring $R$. The following are equivalent
\begin{enumerate}
\item $P$ is a direct summand of a direct sum of modules of the
form $R/fR$, for $f \in R$ varying.
\item for every short exact sequence $0 \to A \to B \to C \to 0$
of $R$-modules such that $fA = A \cap fB$ for all $f \in R$
the map $\Hom_R(P, B) \to \Hom_R(P, C)$ is surjective.
\end{enumerate}
\end{lemma}

\begin{proof}
Let $0 \to A \to B \to C \to 0$ be an exact sequence as in (2).
To prove that (1) implies (2) it suffices to prove that
$\Hom_R(R/fR, B) \to \Hom_R(R/fR, C)$ is surjective for every $f \in R$.
Let $\psi : R/fR \to C$ be a map. Say $\psi(1)$ is the image
of $b \in B$. Then $fb \in A$. Hence there exists an $a \in A$
such that $fa = fb$. Then $f(b - a) = 0$ hence we get a morphism
$\varphi : R/fR \to B$ mapping $1$ to $b - a$ which lifts $\psi$.

\medskip\noindent
Conversely, assume that (2) holds. Let $I$ be the set of pairs
$(f, \varphi)$ where $f \in R$ and $\varphi : R/fR \to P$. For
$i \in I$ denote $(f_i, \varphi_i)$ the corresponding pair.
Consider the map
$$
B = \bigoplus\nolimits_{i \in I} R/f_iR \longrightarrow P
$$
which sends the element $r$ in the summand $R/f_iR$ to $\varphi_i(r)$ in $P$.
Let $A = \Ker(B \to P)$. Then we see that (1) is true if the sequence
$$
0 \to A \to B \to P \to 0
$$
is an exact sequence as in (2). To see this suppose $f \in R$ and
$a \in A$ maps to $f b$ in $B$. Write $b = (r_i)_{i \in I}$ with
almost all $r_i = 0$. Then we see that
$$
f\sum \varphi_i(r_i) = 0
$$
in $P$. Hence there is an $i_0 \in I$ such that $f_{i_0} = f$ and
$\varphi_{i_0}(1) = \sum \varphi_i(r_i)$. Let $x_{i_0} \in R/f_{i_0}R$
be the class of $1$. Then we see that
$$
a' = (r_i)_{i \in I} - (0, \ldots, 0, x_{i_0}, 0, \ldots )
$$
is an element of $A$ and $fa' = a$ as desired.
\end{proof}

\end{document}
